\documentclass[12pt]{article}
\usepackage[utf8]{inputenc}
\usepackage[T1]{fontenc}
\usepackage{hyperref}
\title{Ruwila Language Resources for K-12 Curriculum Development}
\author{Compiled by Claude Sonnet 5.5 (claude-sonnet-5-5) for Andr\'as Kornai}
\date{}
\begin{document}
\maketitle

\section{Introduction}
Ruwila (ISO 639-3 \texttt{rwl}; also Rwila, Lwila, Luwira) is a Bantu language of the Katavi Region of Tanzania. Wikipedia treats it together with Konongo (\texttt{kcz}) as a cluster closely related to Nyamwezi, and some sources count it as a dialect of Nyamwezi. Joshua Project places it in a different Bantu subgroup, so the classification is unsettled. Global Recordings Network and Wikipedia both give about 5,600 speakers (2012). Very little is documented, and I found no school materials in Ruwila.

\section{Primary Educational Resources}
I found no primers, textbooks or school programmes in Ruwila.

\section{Linguistic and Reference Materials}
\begin{itemize}
\item Wikipedia, Konongo--Ruwila language (overview and classification): \url{https://en.wikipedia.org/wiki/Konongo%E2%80%93Ruwila_language}
\item Wikipedia, Malila language (a related Tanzanian language, listed in the same search results, for comparison only): \url{https://en.wikipedia.org/wiki/Malila_language}
\end{itemize}
I found no dictionary, grammar or orthography.

\section{Religious and Community Resources}
\begin{itemize}
\item Joshua Project, Ruwila in Tanzania (people-group profile; location in Mlele District and Nsimbo ward): \url{https://joshuaproject.net/people_groups/21998}
\end{itemize}

\section{Digital and Audio Resources}
\begin{itemize}
\item Global Recordings Network, language page for \texttt{rwl}. It says it has no audio recordings in this language at present, although it lists a Jesus Film version: \url{https://www.globalrecordings.net/en/language/rwl}
\end{itemize}

\section{Curriculum Development Context}
Ruwila has a small speaker base and almost no published documentation. A native-language curriculum would first need community-led documentation: an agreed orthography, a word list and story collection, and the relationship to Nyamwezi and Swahili settled. Swahili is Tanzania's school language.

\end{document}